\section{Conclusion}
\label{sec:conclusion}

\subsection{Summary of contributions}
The MMAudio paper \citep{cheng2024mmaudio} frames V2A as
conditional flow matching with three modalities funnelled through
a single MMDiT stack, and reports state-of-the-art results using
only the flow-matching loss --- no auxiliary contrastive,
alignment, or distribution-matching term. Every cross-modal
coupling it uses (joint attention, Conditional Synchronization
Module, empty-token masking) is pointwise. This thesis adds one
relational-geometry term to that objective and asks, across five
design axes, whether it closes the pointwise-to-relational gap
identified in Sec.\,\ref{sec:intro:pointwise}. The delta we
introduce is exactly one additive loss term; no architectural
change, no pretraining, no new data. Concretely:
\begin{itemize}[leftmargin=*,topsep=2pt]
  \item A batch-level entropic Gromov--Wasserstein regularizer
    $\LGW$ (plus an FGW variant) between intra-video and
    intra-audio pairwise squared-distance matrices, solved by
    Peyr\'e-style mirror descent with an inner Sinkhorn
    projection in float32
    (\texttt{mmaudio/model/gw\_regularization.py}).
  \item Five design axes D1--D5 covering which representations to
    compare, penalty strength, detach direction, schedule, and
    interaction with Synchformer; each owned by one Slurm job
    script under \texttt{jobs/}.
  \item Evaluation beyond aggregate metrics: per-class IB-score
    deltas, representation-geometry diagnostics (effective rank,
    singular spectrum), and coupling-matrix $\Tstar$ heatmaps
    over VGGSound classes.
\end{itemize}

\subsection{Positioning relative to MMAudio's own limitations}
\citet{cheng2024mmaudio} note two scope limitations of MMAudio
in their Section 4.4: (i) the model fails on human speech
because it is specialized to Foley; (ii) scaling past 157\,M
parameters gives diminishing returns, which the authors
attribute to the scarcity of audio-visual data. Our contribution
is orthogonal to both. Relational geometry is a loss-level
property, not a capacity or data-scale property; if H1 holds, a
\emph{fixed-capacity, fixed-data} MMAudio can be improved by
better using the data it already sees. A natural reading of our
hypothesis is: some of the diminishing-return plateau MMAudio
describes is a plateau of pointwise-only supervision, not of
model capacity.

\subsection{Scope and limitations}
\begin{description}[leftmargin=2em]
  \item[Batch-level, not stream-level.] We regularize over the
    mini-batch axis, so $\DV$ and $\DA$ have shape $B \times B$.
    A different reading of ``relational'' would be to regularize
    within a single clip across time (pairwise over the $N=250$
    latent steps). The current design leaves that untouched. A
    time-axis variant would interact closely with Synchformer and
    is a natural extension.
  \item[Per-rank GW under DDP.] As noted in
    Sec.\,\ref{sec:method:subset}, the GW loss is computed on the
    local real-video subset of each rank, not on the global batch.
    This is a consistency approximation; an all-gather before GW
    would give the global-batch objective at the cost of one
    extra NCCL round per step.
  \item[Small\_16k only.] All experiments run at the
    \texttt{small\_16k} configuration; MMAudio also ships
    \texttt{small\_44k}, \texttt{medium\_44k}, and
    \texttt{large\_44k} variants. The training cost ($\sim\!48$
    hours per 300k steps on one GPU) plus the 1 $+$ 20 run grid
    already saturates the cluster budget; scaling to
    \texttt{large\_44k} is future work.
  \item[No training-from-scratch comparisons.] All runs share
    MMAudio's random init but no pretrained base weights are
    loaded. We do not compare against a two-stage fine-tuning
    recipe (pretrain without GW, then fine-tune with GW).
  \item[Sinkhorn numerics are brittle in bfloat16.] This is
    documented in the repository's CLAUDE.md; we work around it
    by forcing float32 inside the GW block. A Sinkhorn-with-log
    formulation would let us reclaim bfloat16 but was not
    necessary for the small batch sizes studied here.
  \item[Metric limitations.] FD$_{\text{PaSST}}$ / IS are
    audio-only distributional; IB-score is one external retrieval
    model; DeSync relies on Synchformer. All four are proxies for
    listener judgment. We do not report human evaluations.
\end{description}

\subsection{Expected contributions}
The contribution we claim is narrow and testable:
\begin{itemize}[leftmargin=*,topsep=2pt]
  \item If H1 holds, entropic GW is a near-free addition to an
    MMAudio-style V2A training objective --- one loss term, one
    scalar, no architecture changes.
  \item If H2 holds, the \emph{level} of representation GW is
    applied to is a first-order design choice, comparable in
    effect to the choice of CLIP vs audio-side projection for a
    contrastive auxiliary.
  \item If H3 holds, the regularizer has a clean operating
    regime (non-monotone in $\lambda_{\mathrm{GW}}$), making it
    tunable rather than brittle.
  \item The coupling-matrix diagnostic in
    \texttt{analyze\_couplings.py} gives a cheap, visual answer
    to ``is relational structure actually captured?'' ---
    independent of whether FD$_{\text{PaSST}}$ went down.
\end{itemize}

\subsection{Future work}
Four extensions follow directly from the limitations:
\begin{enumerate}[leftmargin=*]
  \item Time-axis GW across the $N$ latent steps within a clip,
    possibly replacing or augmenting Synchformer.
  \item Global-batch GW via an all-gather, to quantify the
    per-rank approximation.
  \item GW as a fine-tuning objective on top of a pretrained
    MMAudio checkpoint, to isolate the relational contribution
    from the full 300k-step joint training.
  \item Partial / unbalanced GW \citep{chapel2020partial}:
    relaxing the uniform-marginal assumption would let GW cope
    gracefully with batches that mix genuinely different class
    distributions (e.g. when the oversample ratio changes).
\end{enumerate}
